% sec:ad-comparison — Comparison: AD, Finite Differences, and Analytic Derivatives
\section{Comparison: AD, Finite Differences, and Analytic Derivatives}
\label{sec:ad-comparison}

The geodesic integrator consumes 40 metric derivatives at
every step.  AD produces them exactly --- but a two-line
finite-difference routine can approximate them with no
modification to the metric function at all.  Is the
dual-number machinery of the previous two sections worth the
trouble?  The answer turns on a fundamental accuracy ceiling
that finite differences cannot breach, no matter how
carefully the step size is chosen.

\paragraph{The step-size dilemma.}

Central differences approximate a partial derivative by
evaluating the metric at two perturbed points:
%
\begin{equation}
  \frac{g_{\alpha\beta}(x + h\,e_\mu)
      - g_{\alpha\beta}(x - h\,e_\mu)}{2h}
    \;=\; \pd{\mu} g_{\alpha\beta}
    \;+\; \frac{h^2}{6}\,
      \frac{\partial^3 g_{\alpha\beta}}{\partial(x^\mu)^3}
    \;+\; \order{h^4} \,.
  \label{eq:fd-central-metric}
\end{equation}
%
The truncation error is $\order{h^2}$, just as for the
centred stencils of \cref{sec:fd-truncation}: smaller $h$
means smaller error.  But this argument ignores
floating-point arithmetic.  When $h$ is small enough that
the two metric values agree to many significant digits,
their difference loses those digits to cancellation.  The
result is a fundamental tension: large $h$ introduces
truncation error; small $h$ amplifies round-off.

\paragraph{The accuracy ceiling.}

Quantifying both error sources makes the tradeoff precise.
The truncation term scales as $h^2$ from
\cref{eq:fd-central-metric}.  The round-off term arises
because each computed value $\hat{f}(x \pm h)$ differs from
the true value by at most a relative error of order
$\epsilon_{\text{m}}$, the unit round-off of IEEE 754
double-precision arithmetic
($\epsilon_{\text{m}} = 2^{-52} \approx 2.2 \times 10^{-16}$).
The subtraction in the numerator compounds the error from
both evaluations, and dividing by $2h$ amplifies the result,
so the total error is bounded by
%
\begin{equation}
  E(h) \;\lesssim\;
    \underbrace{\frac{h^2}{6}\,\abs{f'''}}_{\text{truncation}}
    \;+\;
    \underbrace{\frac{\epsilon_{\text{m}}\,\abs{f}}{h}}_{\text{round-off}} \,,
  \label{eq:fd-total-error}
\end{equation}
%
where $f$ stands for any metric component along the
perturbation direction and primes denote ordinary
derivatives with respect to that coordinate.  The two terms
pull in opposite directions: shrinking $h$ reduces the
truncation error but inflates the round-off term.  The best
one can do is balance them.

Minimising \cref{eq:fd-total-error} over $h$ gives the
optimal step size and the best achievable error:
%
\begin{equation}
  h^*
    = \Bigl(
        \frac{3\,\epsilon_{\text{m}}\,\abs{f}}{\abs{f'''}}
      \Bigr)^{\!1/3} ,
  \qquad
  E^* \;\sim\; \epsilon_{\text{m}}^{2/3}\,
    \abs{f}^{2/3}\,\abs{f'''}^{1/3} \,.
  \label{eq:fd-optimal}
\end{equation}
%
The decisive quantity is the exponent: $2/3$, not $1$.  For
double precision,
$\epsilon_{\text{m}}^{2/3} \approx 3.6 \times 10^{-11}$ ---
five orders of magnitude above machine precision.  No choice
of step size can breach this ceiling; it is a fundamental
consequence of approximating derivatives by finite
ratios of function values.
\warnnote{Higher-order stencils (four-point, six-point) raise
the truncation exponent but do not escape the fundamental
limit.  A $p$-th order scheme has truncation error
$\order{h^p}$ and round-off $\order{\epsilon_{\text{m}}/h}$;
balancing gives
$\epsilon_{\text{m}}^{p/(p+1)}$ at best.  Even an
eighth-order stencil reaches only $\sim\!10^{-14}$.}

AD faces no such ceiling.  Every operation in a dual-number
evaluation is individually rounded, but the derivative
propagation is exact at each step, so the accumulated error
is comparable to the rounding of the metric function
itself --- typically a few multiples of
$\epsilon_{\text{m}}$.  The gap between
$\epsilon_{\text{m}}$ and $\epsilon_{\text{m}}^{2/3}$ ---
roughly $10^{-16}$ versus $10^{-11}$ --- is the fundamental
accuracy advantage of AD over finite differences.

The codebase's finite-difference fallback uses a step size
of $h = \max(10^{-6},\; 10^{-6}\,\abs{x^\mu})$, which is
close to the optimal $h^* \sim \epsilon_{\text{m}}^{1/3}
\approx 6 \times 10^{-6}$ --- a well-chosen default, but
one that still leaves the five-order-of-magnitude gap
intact.

\begin{figure}[t]
\centering
%% Creator: Matplotlib, PGF backend
%%
%% To include the figure in your LaTeX document, write
%%   \input{<filename>.pgf}
%%
%% Make sure the required packages are loaded in your preamble
%%   \usepackage{pgf}
%%
%% Also ensure that all the required font packages are loaded; for instance,
%% the lmodern package is sometimes necessary when using math font.
%%   \usepackage{lmodern}
%%
%% Figures using additional raster images can only be included by \input if
%% they are in the same directory as the main LaTeX file. For loading figures
%% from other directories you can use the `import` package
%%   \usepackage{import}
%%
%% and then include the figures with
%%   \import{<path to file>}{<filename>.pgf}
%%
%% Matplotlib used the following preamble
%%   \def\mathdefault#1{#1}
%%   \everymath=\expandafter{\the\everymath\displaystyle}
%%   \IfFileExists{scrextend.sty}{
%%     \usepackage[fontsize=10.000000pt]{scrextend}
%%   }{
%%     \renewcommand{\normalsize}{\fontsize{10.000000}{12.000000}\selectfont}
%%     \normalsize
%%   }
%%   \usepackage{fontspec}
%%   \usepackage{unicode-math}
%%   \setmainfont{STIX Two Text}[Numbers=OldStyle, Ligatures=TeX]
%%   \setmathfont{STIX Two Math}
%%   \setmonofont{Source Code Pro}[Scale=MatchLowercase]
%%   \ifdefined\pdftexversion\else  % non-pdftex case.
%%     \usepackage{fontspec}
%%     \setmainfont{DejaVuSerif.ttf}[Path=\detokenize{/Users/gvn/.pyenv/versions/3.12.11/lib/python3.12/site-packages/matplotlib/mpl-data/fonts/ttf/}]
%%     \setsansfont{DejaVuSans.ttf}[Path=\detokenize{/Users/gvn/.pyenv/versions/3.12.11/lib/python3.12/site-packages/matplotlib/mpl-data/fonts/ttf/}]
%%     \setmonofont{DejaVuSansMono.ttf}[Path=\detokenize{/Users/gvn/.pyenv/versions/3.12.11/lib/python3.12/site-packages/matplotlib/mpl-data/fonts/ttf/}]
%%   \fi
%%   \makeatletter\@ifpackageloaded{underscore}{}{\usepackage[strings]{underscore}}\makeatother
%%
\begingroup%
\makeatletter%
\begin{pgfpicture}%
\pgfpathrectangle{\pgfpointorigin}{\pgfqpoint{3.873422in}{2.319229in}}%
\pgfusepath{use as bounding box, clip}%
\begin{pgfscope}%
\pgfsetbuttcap%
\pgfsetmiterjoin%
\definecolor{currentfill}{rgb}{1.000000,1.000000,1.000000}%
\pgfsetfillcolor{currentfill}%
\pgfsetlinewidth{0.000000pt}%
\definecolor{currentstroke}{rgb}{1.000000,1.000000,1.000000}%
\pgfsetstrokecolor{currentstroke}%
\pgfsetdash{}{0pt}%
\pgfpathmoveto{\pgfqpoint{0.000000in}{0.000000in}}%
\pgfpathlineto{\pgfqpoint{3.873422in}{0.000000in}}%
\pgfpathlineto{\pgfqpoint{3.873422in}{2.319229in}}%
\pgfpathlineto{\pgfqpoint{0.000000in}{2.319229in}}%
\pgfpathlineto{\pgfqpoint{0.000000in}{0.000000in}}%
\pgfpathclose%
\pgfusepath{fill}%
\end{pgfscope}%
\begin{pgfscope}%
\pgfsetbuttcap%
\pgfsetmiterjoin%
\definecolor{currentfill}{rgb}{1.000000,1.000000,1.000000}%
\pgfsetfillcolor{currentfill}%
\pgfsetlinewidth{0.000000pt}%
\definecolor{currentstroke}{rgb}{0.000000,0.000000,0.000000}%
\pgfsetstrokecolor{currentstroke}%
\pgfsetstrokeopacity{0.000000}%
\pgfsetdash{}{0pt}%
\pgfpathmoveto{\pgfqpoint{0.536040in}{0.352669in}}%
\pgfpathlineto{\pgfqpoint{3.708890in}{0.352669in}}%
\pgfpathlineto{\pgfqpoint{3.708890in}{2.299229in}}%
\pgfpathlineto{\pgfqpoint{0.536040in}{2.299229in}}%
\pgfpathlineto{\pgfqpoint{0.536040in}{0.352669in}}%
\pgfpathclose%
\pgfusepath{fill}%
\end{pgfscope}%
\begin{pgfscope}%
\pgfpathrectangle{\pgfqpoint{0.536040in}{0.352669in}}{\pgfqpoint{3.172850in}{1.946560in}}%
\pgfusepath{clip}%
\pgfsetbuttcap%
\pgfsetroundjoin%
\pgfsetlinewidth{0.702625pt}%
\definecolor{currentstroke}{rgb}{0.431373,0.423529,0.462745}%
\pgfsetstrokecolor{currentstroke}%
\pgfsetstrokeopacity{0.600000}%
\pgfsetdash{{0.700000pt}{1.155000pt}}{0.000000pt}%
\pgfpathmoveto{\pgfqpoint{2.016703in}{0.356496in}}%
\pgfpathlineto{\pgfqpoint{2.032728in}{0.373845in}}%
\pgfpathlineto{\pgfqpoint{2.048752in}{0.391194in}}%
\pgfpathlineto{\pgfqpoint{2.064777in}{0.408543in}}%
\pgfpathlineto{\pgfqpoint{2.080801in}{0.425892in}}%
\pgfpathlineto{\pgfqpoint{2.096826in}{0.443241in}}%
\pgfpathlineto{\pgfqpoint{2.112850in}{0.460590in}}%
\pgfpathlineto{\pgfqpoint{2.128875in}{0.477939in}}%
\pgfpathlineto{\pgfqpoint{2.144899in}{0.495288in}}%
\pgfpathlineto{\pgfqpoint{2.160924in}{0.512637in}}%
\pgfpathlineto{\pgfqpoint{2.176948in}{0.529986in}}%
\pgfpathlineto{\pgfqpoint{2.192973in}{0.547335in}}%
\pgfpathlineto{\pgfqpoint{2.208997in}{0.564684in}}%
\pgfpathlineto{\pgfqpoint{2.225022in}{0.582033in}}%
\pgfpathlineto{\pgfqpoint{2.241046in}{0.599382in}}%
\pgfpathlineto{\pgfqpoint{2.257071in}{0.616731in}}%
\pgfpathlineto{\pgfqpoint{2.273095in}{0.634080in}}%
\pgfpathlineto{\pgfqpoint{2.289120in}{0.651429in}}%
\pgfpathlineto{\pgfqpoint{2.305144in}{0.668778in}}%
\pgfpathlineto{\pgfqpoint{2.321169in}{0.686127in}}%
\pgfpathlineto{\pgfqpoint{2.337193in}{0.703476in}}%
\pgfpathlineto{\pgfqpoint{2.353218in}{0.720825in}}%
\pgfpathlineto{\pgfqpoint{2.369242in}{0.738174in}}%
\pgfpathlineto{\pgfqpoint{2.385267in}{0.755523in}}%
\pgfpathlineto{\pgfqpoint{2.401291in}{0.772872in}}%
\pgfpathlineto{\pgfqpoint{2.417316in}{0.790221in}}%
\pgfpathlineto{\pgfqpoint{2.433340in}{0.807570in}}%
\pgfpathlineto{\pgfqpoint{2.449365in}{0.824919in}}%
\pgfpathlineto{\pgfqpoint{2.465389in}{0.842268in}}%
\pgfpathlineto{\pgfqpoint{2.481414in}{0.859617in}}%
\pgfpathlineto{\pgfqpoint{2.497438in}{0.876966in}}%
\pgfpathlineto{\pgfqpoint{2.513463in}{0.894315in}}%
\pgfpathlineto{\pgfqpoint{2.529487in}{0.911665in}}%
\pgfpathlineto{\pgfqpoint{2.545512in}{0.929014in}}%
\pgfpathlineto{\pgfqpoint{2.561536in}{0.946363in}}%
\pgfpathlineto{\pgfqpoint{2.577561in}{0.963712in}}%
\pgfpathlineto{\pgfqpoint{2.593585in}{0.981061in}}%
\pgfpathlineto{\pgfqpoint{2.609610in}{0.998410in}}%
\pgfpathlineto{\pgfqpoint{2.625634in}{1.015759in}}%
\pgfpathlineto{\pgfqpoint{2.641659in}{1.033108in}}%
\pgfpathlineto{\pgfqpoint{2.657683in}{1.050457in}}%
\pgfpathlineto{\pgfqpoint{2.673708in}{1.067806in}}%
\pgfpathlineto{\pgfqpoint{2.689732in}{1.085155in}}%
\pgfpathlineto{\pgfqpoint{2.705757in}{1.102504in}}%
\pgfpathlineto{\pgfqpoint{2.721781in}{1.119853in}}%
\pgfpathlineto{\pgfqpoint{2.737806in}{1.137202in}}%
\pgfpathlineto{\pgfqpoint{2.753830in}{1.154551in}}%
\pgfpathlineto{\pgfqpoint{2.769855in}{1.171900in}}%
\pgfpathlineto{\pgfqpoint{2.785879in}{1.189249in}}%
\pgfpathlineto{\pgfqpoint{2.801904in}{1.206598in}}%
\pgfpathlineto{\pgfqpoint{2.817928in}{1.223947in}}%
\pgfpathlineto{\pgfqpoint{2.833953in}{1.241296in}}%
\pgfpathlineto{\pgfqpoint{2.849977in}{1.258645in}}%
\pgfpathlineto{\pgfqpoint{2.866002in}{1.275994in}}%
\pgfpathlineto{\pgfqpoint{2.882026in}{1.293343in}}%
\pgfpathlineto{\pgfqpoint{2.898051in}{1.310692in}}%
\pgfpathlineto{\pgfqpoint{2.914075in}{1.328041in}}%
\pgfpathlineto{\pgfqpoint{2.930100in}{1.345390in}}%
\pgfpathlineto{\pgfqpoint{2.946124in}{1.362739in}}%
\pgfpathlineto{\pgfqpoint{2.962149in}{1.380088in}}%
\pgfpathlineto{\pgfqpoint{2.978173in}{1.397437in}}%
\pgfpathlineto{\pgfqpoint{2.994198in}{1.414786in}}%
\pgfpathlineto{\pgfqpoint{3.010222in}{1.432135in}}%
\pgfpathlineto{\pgfqpoint{3.026247in}{1.449484in}}%
\pgfpathlineto{\pgfqpoint{3.042271in}{1.466833in}}%
\pgfpathlineto{\pgfqpoint{3.058296in}{1.484182in}}%
\pgfpathlineto{\pgfqpoint{3.074320in}{1.501531in}}%
\pgfpathlineto{\pgfqpoint{3.090345in}{1.518880in}}%
\pgfpathlineto{\pgfqpoint{3.106369in}{1.536229in}}%
\pgfpathlineto{\pgfqpoint{3.122394in}{1.553578in}}%
\pgfpathlineto{\pgfqpoint{3.138418in}{1.570927in}}%
\pgfpathlineto{\pgfqpoint{3.154443in}{1.588276in}}%
\pgfpathlineto{\pgfqpoint{3.170467in}{1.605625in}}%
\pgfpathlineto{\pgfqpoint{3.186492in}{1.622974in}}%
\pgfpathlineto{\pgfqpoint{3.202516in}{1.640323in}}%
\pgfpathlineto{\pgfqpoint{3.218541in}{1.657672in}}%
\pgfpathlineto{\pgfqpoint{3.234565in}{1.675021in}}%
\pgfpathlineto{\pgfqpoint{3.250590in}{1.692370in}}%
\pgfpathlineto{\pgfqpoint{3.266614in}{1.709719in}}%
\pgfpathlineto{\pgfqpoint{3.282639in}{1.727068in}}%
\pgfpathlineto{\pgfqpoint{3.298663in}{1.744417in}}%
\pgfpathlineto{\pgfqpoint{3.314688in}{1.761766in}}%
\pgfpathlineto{\pgfqpoint{3.330712in}{1.779115in}}%
\pgfpathlineto{\pgfqpoint{3.346737in}{1.796465in}}%
\pgfpathlineto{\pgfqpoint{3.362761in}{1.813814in}}%
\pgfpathlineto{\pgfqpoint{3.378786in}{1.831163in}}%
\pgfpathlineto{\pgfqpoint{3.394810in}{1.848512in}}%
\pgfpathlineto{\pgfqpoint{3.410835in}{1.865861in}}%
\pgfpathlineto{\pgfqpoint{3.426859in}{1.883210in}}%
\pgfpathlineto{\pgfqpoint{3.442884in}{1.900559in}}%
\pgfpathlineto{\pgfqpoint{3.458908in}{1.917908in}}%
\pgfpathlineto{\pgfqpoint{3.474933in}{1.935257in}}%
\pgfpathlineto{\pgfqpoint{3.490957in}{1.952606in}}%
\pgfpathlineto{\pgfqpoint{3.506982in}{1.969955in}}%
\pgfpathlineto{\pgfqpoint{3.523006in}{1.987304in}}%
\pgfpathlineto{\pgfqpoint{3.539030in}{2.004653in}}%
\pgfpathlineto{\pgfqpoint{3.555055in}{2.022002in}}%
\pgfpathlineto{\pgfqpoint{3.571079in}{2.039351in}}%
\pgfpathlineto{\pgfqpoint{3.587104in}{2.056700in}}%
\pgfpathlineto{\pgfqpoint{3.603128in}{2.074049in}}%
\pgfusepath{stroke}%
\end{pgfscope}%
\begin{pgfscope}%
\pgfpathrectangle{\pgfqpoint{0.536040in}{0.352669in}}{\pgfqpoint{3.172850in}{1.946560in}}%
\pgfusepath{clip}%
\pgfsetbuttcap%
\pgfsetroundjoin%
\pgfsetlinewidth{0.702625pt}%
\definecolor{currentstroke}{rgb}{0.431373,0.423529,0.462745}%
\pgfsetstrokecolor{currentstroke}%
\pgfsetstrokeopacity{0.600000}%
\pgfsetdash{{0.700000pt}{1.155000pt}}{0.000000pt}%
\pgfpathmoveto{\pgfqpoint{2.016703in}{1.446899in}}%
\pgfpathlineto{\pgfqpoint{2.032728in}{1.438225in}}%
\pgfpathlineto{\pgfqpoint{2.048752in}{1.429550in}}%
\pgfpathlineto{\pgfqpoint{2.064777in}{1.420876in}}%
\pgfpathlineto{\pgfqpoint{2.080801in}{1.412201in}}%
\pgfpathlineto{\pgfqpoint{2.096826in}{1.403527in}}%
\pgfpathlineto{\pgfqpoint{2.112850in}{1.394852in}}%
\pgfpathlineto{\pgfqpoint{2.128875in}{1.386178in}}%
\pgfpathlineto{\pgfqpoint{2.144899in}{1.377503in}}%
\pgfpathlineto{\pgfqpoint{2.160924in}{1.368829in}}%
\pgfpathlineto{\pgfqpoint{2.176948in}{1.360154in}}%
\pgfpathlineto{\pgfqpoint{2.192973in}{1.351480in}}%
\pgfpathlineto{\pgfqpoint{2.208997in}{1.342805in}}%
\pgfpathlineto{\pgfqpoint{2.225022in}{1.334130in}}%
\pgfpathlineto{\pgfqpoint{2.241046in}{1.325456in}}%
\pgfpathlineto{\pgfqpoint{2.257071in}{1.316781in}}%
\pgfpathlineto{\pgfqpoint{2.273095in}{1.308107in}}%
\pgfpathlineto{\pgfqpoint{2.289120in}{1.299432in}}%
\pgfpathlineto{\pgfqpoint{2.305144in}{1.290758in}}%
\pgfpathlineto{\pgfqpoint{2.321169in}{1.282083in}}%
\pgfpathlineto{\pgfqpoint{2.337193in}{1.273409in}}%
\pgfpathlineto{\pgfqpoint{2.353218in}{1.264734in}}%
\pgfpathlineto{\pgfqpoint{2.369242in}{1.256060in}}%
\pgfpathlineto{\pgfqpoint{2.385267in}{1.247385in}}%
\pgfpathlineto{\pgfqpoint{2.401291in}{1.238711in}}%
\pgfpathlineto{\pgfqpoint{2.417316in}{1.230036in}}%
\pgfpathlineto{\pgfqpoint{2.433340in}{1.221362in}}%
\pgfpathlineto{\pgfqpoint{2.449365in}{1.212687in}}%
\pgfpathlineto{\pgfqpoint{2.465389in}{1.204013in}}%
\pgfpathlineto{\pgfqpoint{2.481414in}{1.195338in}}%
\pgfpathlineto{\pgfqpoint{2.497438in}{1.186664in}}%
\pgfpathlineto{\pgfqpoint{2.513463in}{1.177989in}}%
\pgfpathlineto{\pgfqpoint{2.529487in}{1.169315in}}%
\pgfpathlineto{\pgfqpoint{2.545512in}{1.160640in}}%
\pgfpathlineto{\pgfqpoint{2.561536in}{1.151966in}}%
\pgfpathlineto{\pgfqpoint{2.577561in}{1.143291in}}%
\pgfpathlineto{\pgfqpoint{2.593585in}{1.134617in}}%
\pgfpathlineto{\pgfqpoint{2.609610in}{1.125942in}}%
\pgfpathlineto{\pgfqpoint{2.625634in}{1.117268in}}%
\pgfpathlineto{\pgfqpoint{2.641659in}{1.108593in}}%
\pgfpathlineto{\pgfqpoint{2.657683in}{1.099919in}}%
\pgfpathlineto{\pgfqpoint{2.673708in}{1.091244in}}%
\pgfpathlineto{\pgfqpoint{2.689732in}{1.082570in}}%
\pgfpathlineto{\pgfqpoint{2.705757in}{1.073895in}}%
\pgfpathlineto{\pgfqpoint{2.721781in}{1.065221in}}%
\pgfpathlineto{\pgfqpoint{2.737806in}{1.056546in}}%
\pgfpathlineto{\pgfqpoint{2.753830in}{1.047872in}}%
\pgfpathlineto{\pgfqpoint{2.769855in}{1.039197in}}%
\pgfpathlineto{\pgfqpoint{2.785879in}{1.030523in}}%
\pgfpathlineto{\pgfqpoint{2.801904in}{1.021848in}}%
\pgfpathlineto{\pgfqpoint{2.817928in}{1.013174in}}%
\pgfpathlineto{\pgfqpoint{2.833953in}{1.004499in}}%
\pgfpathlineto{\pgfqpoint{2.849977in}{0.995825in}}%
\pgfpathlineto{\pgfqpoint{2.866002in}{0.987150in}}%
\pgfpathlineto{\pgfqpoint{2.882026in}{0.978476in}}%
\pgfpathlineto{\pgfqpoint{2.898051in}{0.969801in}}%
\pgfpathlineto{\pgfqpoint{2.914075in}{0.961127in}}%
\pgfpathlineto{\pgfqpoint{2.930100in}{0.952452in}}%
\pgfpathlineto{\pgfqpoint{2.946124in}{0.943778in}}%
\pgfpathlineto{\pgfqpoint{2.962149in}{0.935103in}}%
\pgfpathlineto{\pgfqpoint{2.978173in}{0.926429in}}%
\pgfpathlineto{\pgfqpoint{2.994198in}{0.917754in}}%
\pgfpathlineto{\pgfqpoint{3.010222in}{0.909080in}}%
\pgfpathlineto{\pgfqpoint{3.026247in}{0.900405in}}%
\pgfpathlineto{\pgfqpoint{3.042271in}{0.891730in}}%
\pgfpathlineto{\pgfqpoint{3.058296in}{0.883056in}}%
\pgfpathlineto{\pgfqpoint{3.074320in}{0.874381in}}%
\pgfpathlineto{\pgfqpoint{3.090345in}{0.865707in}}%
\pgfpathlineto{\pgfqpoint{3.106369in}{0.857032in}}%
\pgfpathlineto{\pgfqpoint{3.122394in}{0.848358in}}%
\pgfpathlineto{\pgfqpoint{3.138418in}{0.839683in}}%
\pgfpathlineto{\pgfqpoint{3.154443in}{0.831009in}}%
\pgfpathlineto{\pgfqpoint{3.170467in}{0.822334in}}%
\pgfpathlineto{\pgfqpoint{3.186492in}{0.813660in}}%
\pgfpathlineto{\pgfqpoint{3.202516in}{0.804985in}}%
\pgfpathlineto{\pgfqpoint{3.218541in}{0.796311in}}%
\pgfpathlineto{\pgfqpoint{3.234565in}{0.787636in}}%
\pgfpathlineto{\pgfqpoint{3.250590in}{0.778962in}}%
\pgfpathlineto{\pgfqpoint{3.266614in}{0.770287in}}%
\pgfpathlineto{\pgfqpoint{3.282639in}{0.761613in}}%
\pgfpathlineto{\pgfqpoint{3.298663in}{0.752938in}}%
\pgfpathlineto{\pgfqpoint{3.314688in}{0.744264in}}%
\pgfpathlineto{\pgfqpoint{3.330712in}{0.735589in}}%
\pgfpathlineto{\pgfqpoint{3.346737in}{0.726915in}}%
\pgfpathlineto{\pgfqpoint{3.362761in}{0.718240in}}%
\pgfpathlineto{\pgfqpoint{3.378786in}{0.709566in}}%
\pgfpathlineto{\pgfqpoint{3.394810in}{0.700891in}}%
\pgfpathlineto{\pgfqpoint{3.410835in}{0.692217in}}%
\pgfpathlineto{\pgfqpoint{3.426859in}{0.683542in}}%
\pgfpathlineto{\pgfqpoint{3.442884in}{0.674868in}}%
\pgfpathlineto{\pgfqpoint{3.458908in}{0.666193in}}%
\pgfpathlineto{\pgfqpoint{3.474933in}{0.657519in}}%
\pgfpathlineto{\pgfqpoint{3.490957in}{0.648844in}}%
\pgfpathlineto{\pgfqpoint{3.506982in}{0.640170in}}%
\pgfpathlineto{\pgfqpoint{3.523006in}{0.631495in}}%
\pgfpathlineto{\pgfqpoint{3.539030in}{0.622821in}}%
\pgfpathlineto{\pgfqpoint{3.555055in}{0.614146in}}%
\pgfpathlineto{\pgfqpoint{3.571079in}{0.605472in}}%
\pgfpathlineto{\pgfqpoint{3.587104in}{0.596797in}}%
\pgfpathlineto{\pgfqpoint{3.603128in}{0.588123in}}%
\pgfusepath{stroke}%
\end{pgfscope}%
\begin{pgfscope}%
\pgfsetbuttcap%
\pgfsetroundjoin%
\definecolor{currentfill}{rgb}{0.000000,0.000000,0.000000}%
\pgfsetfillcolor{currentfill}%
\pgfsetlinewidth{0.501875pt}%
\definecolor{currentstroke}{rgb}{0.000000,0.000000,0.000000}%
\pgfsetstrokecolor{currentstroke}%
\pgfsetdash{}{0pt}%
\pgfsys@defobject{currentmarker}{\pgfqpoint{0.000000in}{0.000000in}}{\pgfqpoint{0.000000in}{0.041667in}}{%
\pgfpathmoveto{\pgfqpoint{0.000000in}{0.000000in}}%
\pgfpathlineto{\pgfqpoint{0.000000in}{0.041667in}}%
\pgfusepath{stroke,fill}%
}%
\begin{pgfscope}%
\pgfsys@transformshift{0.536040in}{0.352669in}%
\pgfsys@useobject{currentmarker}{}%
\end{pgfscope}%
\end{pgfscope}%
\begin{pgfscope}%
\definecolor{textcolor}{rgb}{0.000000,0.000000,0.000000}%
\pgfsetstrokecolor{textcolor}%
\pgfsetfillcolor{textcolor}%
\pgftext[x=0.536040in,y=0.304058in,,top]{\color{textcolor}{\sffamily\fontsize{8.000000}{9.600000}\selectfont\catcode`\^=\active\def^{\ifmmode\sp\else\^{}\fi}\catcode`\%=\active\def%{\%}$\mathdefault{10^{-15}}$}}%
\end{pgfscope}%
\begin{pgfscope}%
\pgfsetbuttcap%
\pgfsetroundjoin%
\definecolor{currentfill}{rgb}{0.000000,0.000000,0.000000}%
\pgfsetfillcolor{currentfill}%
\pgfsetlinewidth{0.501875pt}%
\definecolor{currentstroke}{rgb}{0.000000,0.000000,0.000000}%
\pgfsetstrokecolor{currentstroke}%
\pgfsetdash{}{0pt}%
\pgfsys@defobject{currentmarker}{\pgfqpoint{0.000000in}{0.000000in}}{\pgfqpoint{0.000000in}{0.041667in}}{%
\pgfpathmoveto{\pgfqpoint{0.000000in}{0.000000in}}%
\pgfpathlineto{\pgfqpoint{0.000000in}{0.041667in}}%
\pgfusepath{stroke,fill}%
}%
\begin{pgfscope}%
\pgfsys@transformshift{0.959087in}{0.352669in}%
\pgfsys@useobject{currentmarker}{}%
\end{pgfscope}%
\end{pgfscope}%
\begin{pgfscope}%
\definecolor{textcolor}{rgb}{0.000000,0.000000,0.000000}%
\pgfsetstrokecolor{textcolor}%
\pgfsetfillcolor{textcolor}%
\pgftext[x=0.959087in,y=0.304058in,,top]{\color{textcolor}{\sffamily\fontsize{8.000000}{9.600000}\selectfont\catcode`\^=\active\def^{\ifmmode\sp\else\^{}\fi}\catcode`\%=\active\def%{\%}$\mathdefault{10^{-13}}$}}%
\end{pgfscope}%
\begin{pgfscope}%
\pgfsetbuttcap%
\pgfsetroundjoin%
\definecolor{currentfill}{rgb}{0.000000,0.000000,0.000000}%
\pgfsetfillcolor{currentfill}%
\pgfsetlinewidth{0.501875pt}%
\definecolor{currentstroke}{rgb}{0.000000,0.000000,0.000000}%
\pgfsetstrokecolor{currentstroke}%
\pgfsetdash{}{0pt}%
\pgfsys@defobject{currentmarker}{\pgfqpoint{0.000000in}{0.000000in}}{\pgfqpoint{0.000000in}{0.041667in}}{%
\pgfpathmoveto{\pgfqpoint{0.000000in}{0.000000in}}%
\pgfpathlineto{\pgfqpoint{0.000000in}{0.041667in}}%
\pgfusepath{stroke,fill}%
}%
\begin{pgfscope}%
\pgfsys@transformshift{1.382133in}{0.352669in}%
\pgfsys@useobject{currentmarker}{}%
\end{pgfscope}%
\end{pgfscope}%
\begin{pgfscope}%
\definecolor{textcolor}{rgb}{0.000000,0.000000,0.000000}%
\pgfsetstrokecolor{textcolor}%
\pgfsetfillcolor{textcolor}%
\pgftext[x=1.382133in,y=0.304058in,,top]{\color{textcolor}{\sffamily\fontsize{8.000000}{9.600000}\selectfont\catcode`\^=\active\def^{\ifmmode\sp\else\^{}\fi}\catcode`\%=\active\def%{\%}$\mathdefault{10^{-11}}$}}%
\end{pgfscope}%
\begin{pgfscope}%
\pgfsetbuttcap%
\pgfsetroundjoin%
\definecolor{currentfill}{rgb}{0.000000,0.000000,0.000000}%
\pgfsetfillcolor{currentfill}%
\pgfsetlinewidth{0.501875pt}%
\definecolor{currentstroke}{rgb}{0.000000,0.000000,0.000000}%
\pgfsetstrokecolor{currentstroke}%
\pgfsetdash{}{0pt}%
\pgfsys@defobject{currentmarker}{\pgfqpoint{0.000000in}{0.000000in}}{\pgfqpoint{0.000000in}{0.041667in}}{%
\pgfpathmoveto{\pgfqpoint{0.000000in}{0.000000in}}%
\pgfpathlineto{\pgfqpoint{0.000000in}{0.041667in}}%
\pgfusepath{stroke,fill}%
}%
\begin{pgfscope}%
\pgfsys@transformshift{1.805180in}{0.352669in}%
\pgfsys@useobject{currentmarker}{}%
\end{pgfscope}%
\end{pgfscope}%
\begin{pgfscope}%
\definecolor{textcolor}{rgb}{0.000000,0.000000,0.000000}%
\pgfsetstrokecolor{textcolor}%
\pgfsetfillcolor{textcolor}%
\pgftext[x=1.805180in,y=0.304058in,,top]{\color{textcolor}{\sffamily\fontsize{8.000000}{9.600000}\selectfont\catcode`\^=\active\def^{\ifmmode\sp\else\^{}\fi}\catcode`\%=\active\def%{\%}$\mathdefault{10^{-9}}$}}%
\end{pgfscope}%
\begin{pgfscope}%
\pgfsetbuttcap%
\pgfsetroundjoin%
\definecolor{currentfill}{rgb}{0.000000,0.000000,0.000000}%
\pgfsetfillcolor{currentfill}%
\pgfsetlinewidth{0.501875pt}%
\definecolor{currentstroke}{rgb}{0.000000,0.000000,0.000000}%
\pgfsetstrokecolor{currentstroke}%
\pgfsetdash{}{0pt}%
\pgfsys@defobject{currentmarker}{\pgfqpoint{0.000000in}{0.000000in}}{\pgfqpoint{0.000000in}{0.041667in}}{%
\pgfpathmoveto{\pgfqpoint{0.000000in}{0.000000in}}%
\pgfpathlineto{\pgfqpoint{0.000000in}{0.041667in}}%
\pgfusepath{stroke,fill}%
}%
\begin{pgfscope}%
\pgfsys@transformshift{2.228227in}{0.352669in}%
\pgfsys@useobject{currentmarker}{}%
\end{pgfscope}%
\end{pgfscope}%
\begin{pgfscope}%
\definecolor{textcolor}{rgb}{0.000000,0.000000,0.000000}%
\pgfsetstrokecolor{textcolor}%
\pgfsetfillcolor{textcolor}%
\pgftext[x=2.228227in,y=0.304058in,,top]{\color{textcolor}{\sffamily\fontsize{8.000000}{9.600000}\selectfont\catcode`\^=\active\def^{\ifmmode\sp\else\^{}\fi}\catcode`\%=\active\def%{\%}$\mathdefault{10^{-7}}$}}%
\end{pgfscope}%
\begin{pgfscope}%
\pgfsetbuttcap%
\pgfsetroundjoin%
\definecolor{currentfill}{rgb}{0.000000,0.000000,0.000000}%
\pgfsetfillcolor{currentfill}%
\pgfsetlinewidth{0.501875pt}%
\definecolor{currentstroke}{rgb}{0.000000,0.000000,0.000000}%
\pgfsetstrokecolor{currentstroke}%
\pgfsetdash{}{0pt}%
\pgfsys@defobject{currentmarker}{\pgfqpoint{0.000000in}{0.000000in}}{\pgfqpoint{0.000000in}{0.041667in}}{%
\pgfpathmoveto{\pgfqpoint{0.000000in}{0.000000in}}%
\pgfpathlineto{\pgfqpoint{0.000000in}{0.041667in}}%
\pgfusepath{stroke,fill}%
}%
\begin{pgfscope}%
\pgfsys@transformshift{2.651273in}{0.352669in}%
\pgfsys@useobject{currentmarker}{}%
\end{pgfscope}%
\end{pgfscope}%
\begin{pgfscope}%
\definecolor{textcolor}{rgb}{0.000000,0.000000,0.000000}%
\pgfsetstrokecolor{textcolor}%
\pgfsetfillcolor{textcolor}%
\pgftext[x=2.651273in,y=0.304058in,,top]{\color{textcolor}{\sffamily\fontsize{8.000000}{9.600000}\selectfont\catcode`\^=\active\def^{\ifmmode\sp\else\^{}\fi}\catcode`\%=\active\def%{\%}$\mathdefault{10^{-5}}$}}%
\end{pgfscope}%
\begin{pgfscope}%
\pgfsetbuttcap%
\pgfsetroundjoin%
\definecolor{currentfill}{rgb}{0.000000,0.000000,0.000000}%
\pgfsetfillcolor{currentfill}%
\pgfsetlinewidth{0.501875pt}%
\definecolor{currentstroke}{rgb}{0.000000,0.000000,0.000000}%
\pgfsetstrokecolor{currentstroke}%
\pgfsetdash{}{0pt}%
\pgfsys@defobject{currentmarker}{\pgfqpoint{0.000000in}{0.000000in}}{\pgfqpoint{0.000000in}{0.041667in}}{%
\pgfpathmoveto{\pgfqpoint{0.000000in}{0.000000in}}%
\pgfpathlineto{\pgfqpoint{0.000000in}{0.041667in}}%
\pgfusepath{stroke,fill}%
}%
\begin{pgfscope}%
\pgfsys@transformshift{3.074320in}{0.352669in}%
\pgfsys@useobject{currentmarker}{}%
\end{pgfscope}%
\end{pgfscope}%
\begin{pgfscope}%
\definecolor{textcolor}{rgb}{0.000000,0.000000,0.000000}%
\pgfsetstrokecolor{textcolor}%
\pgfsetfillcolor{textcolor}%
\pgftext[x=3.074320in,y=0.304058in,,top]{\color{textcolor}{\sffamily\fontsize{8.000000}{9.600000}\selectfont\catcode`\^=\active\def^{\ifmmode\sp\else\^{}\fi}\catcode`\%=\active\def%{\%}$\mathdefault{10^{-3}}$}}%
\end{pgfscope}%
\begin{pgfscope}%
\pgfsetbuttcap%
\pgfsetroundjoin%
\definecolor{currentfill}{rgb}{0.000000,0.000000,0.000000}%
\pgfsetfillcolor{currentfill}%
\pgfsetlinewidth{0.501875pt}%
\definecolor{currentstroke}{rgb}{0.000000,0.000000,0.000000}%
\pgfsetstrokecolor{currentstroke}%
\pgfsetdash{}{0pt}%
\pgfsys@defobject{currentmarker}{\pgfqpoint{0.000000in}{0.000000in}}{\pgfqpoint{0.000000in}{0.041667in}}{%
\pgfpathmoveto{\pgfqpoint{0.000000in}{0.000000in}}%
\pgfpathlineto{\pgfqpoint{0.000000in}{0.041667in}}%
\pgfusepath{stroke,fill}%
}%
\begin{pgfscope}%
\pgfsys@transformshift{3.497367in}{0.352669in}%
\pgfsys@useobject{currentmarker}{}%
\end{pgfscope}%
\end{pgfscope}%
\begin{pgfscope}%
\definecolor{textcolor}{rgb}{0.000000,0.000000,0.000000}%
\pgfsetstrokecolor{textcolor}%
\pgfsetfillcolor{textcolor}%
\pgftext[x=3.497367in,y=0.304058in,,top]{\color{textcolor}{\sffamily\fontsize{8.000000}{9.600000}\selectfont\catcode`\^=\active\def^{\ifmmode\sp\else\^{}\fi}\catcode`\%=\active\def%{\%}$\mathdefault{10^{-1}}$}}%
\end{pgfscope}%
\begin{pgfscope}%
\definecolor{textcolor}{rgb}{0.000000,0.000000,0.000000}%
\pgfsetstrokecolor{textcolor}%
\pgfsetfillcolor{textcolor}%
\pgftext[x=2.122465in,y=0.140972in,,top]{\color{textcolor}{\sffamily\fontsize{9.000000}{10.800000}\selectfont\catcode`\^=\active\def^{\ifmmode\sp\else\^{}\fi}\catcode`\%=\active\def%{\%}FD step size $h$}}%
\end{pgfscope}%
\begin{pgfscope}%
\pgfsetbuttcap%
\pgfsetroundjoin%
\definecolor{currentfill}{rgb}{0.000000,0.000000,0.000000}%
\pgfsetfillcolor{currentfill}%
\pgfsetlinewidth{0.501875pt}%
\definecolor{currentstroke}{rgb}{0.000000,0.000000,0.000000}%
\pgfsetstrokecolor{currentstroke}%
\pgfsetdash{}{0pt}%
\pgfsys@defobject{currentmarker}{\pgfqpoint{0.000000in}{0.000000in}}{\pgfqpoint{0.041667in}{0.000000in}}{%
\pgfpathmoveto{\pgfqpoint{0.000000in}{0.000000in}}%
\pgfpathlineto{\pgfqpoint{0.041667in}{0.000000in}}%
\pgfusepath{stroke,fill}%
}%
\begin{pgfscope}%
\pgfsys@transformshift{0.536040in}{0.352669in}%
\pgfsys@useobject{currentmarker}{}%
\end{pgfscope}%
\end{pgfscope}%
\begin{pgfscope}%
\definecolor{textcolor}{rgb}{0.000000,0.000000,0.000000}%
\pgfsetstrokecolor{textcolor}%
\pgfsetfillcolor{textcolor}%
\pgftext[x=0.203734in, y=0.310460in, left, base]{\color{textcolor}{\sffamily\fontsize{8.000000}{9.600000}\selectfont\catcode`\^=\active\def^{\ifmmode\sp\else\^{}\fi}\catcode`\%=\active\def%{\%}$\mathdefault{10^{-17}}$}}%
\end{pgfscope}%
\begin{pgfscope}%
\pgfsetbuttcap%
\pgfsetroundjoin%
\definecolor{currentfill}{rgb}{0.000000,0.000000,0.000000}%
\pgfsetfillcolor{currentfill}%
\pgfsetlinewidth{0.501875pt}%
\definecolor{currentstroke}{rgb}{0.000000,0.000000,0.000000}%
\pgfsetstrokecolor{currentstroke}%
\pgfsetdash{}{0pt}%
\pgfsys@defobject{currentmarker}{\pgfqpoint{0.000000in}{0.000000in}}{\pgfqpoint{0.041667in}{0.000000in}}{%
\pgfpathmoveto{\pgfqpoint{0.000000in}{0.000000in}}%
\pgfpathlineto{\pgfqpoint{0.041667in}{0.000000in}}%
\pgfusepath{stroke,fill}%
}%
\begin{pgfscope}%
\pgfsys@transformshift{0.536040in}{0.696179in}%
\pgfsys@useobject{currentmarker}{}%
\end{pgfscope}%
\end{pgfscope}%
\begin{pgfscope}%
\definecolor{textcolor}{rgb}{0.000000,0.000000,0.000000}%
\pgfsetstrokecolor{textcolor}%
\pgfsetfillcolor{textcolor}%
\pgftext[x=0.198651in, y=0.653970in, left, base]{\color{textcolor}{\sffamily\fontsize{8.000000}{9.600000}\selectfont\catcode`\^=\active\def^{\ifmmode\sp\else\^{}\fi}\catcode`\%=\active\def%{\%}$\mathdefault{10^{-14}}$}}%
\end{pgfscope}%
\begin{pgfscope}%
\pgfsetbuttcap%
\pgfsetroundjoin%
\definecolor{currentfill}{rgb}{0.000000,0.000000,0.000000}%
\pgfsetfillcolor{currentfill}%
\pgfsetlinewidth{0.501875pt}%
\definecolor{currentstroke}{rgb}{0.000000,0.000000,0.000000}%
\pgfsetstrokecolor{currentstroke}%
\pgfsetdash{}{0pt}%
\pgfsys@defobject{currentmarker}{\pgfqpoint{0.000000in}{0.000000in}}{\pgfqpoint{0.041667in}{0.000000in}}{%
\pgfpathmoveto{\pgfqpoint{0.000000in}{0.000000in}}%
\pgfpathlineto{\pgfqpoint{0.041667in}{0.000000in}}%
\pgfusepath{stroke,fill}%
}%
\begin{pgfscope}%
\pgfsys@transformshift{0.536040in}{1.039690in}%
\pgfsys@useobject{currentmarker}{}%
\end{pgfscope}%
\end{pgfscope}%
\begin{pgfscope}%
\definecolor{textcolor}{rgb}{0.000000,0.000000,0.000000}%
\pgfsetstrokecolor{textcolor}%
\pgfsetfillcolor{textcolor}%
\pgftext[x=0.201818in, y=0.997481in, left, base]{\color{textcolor}{\sffamily\fontsize{8.000000}{9.600000}\selectfont\catcode`\^=\active\def^{\ifmmode\sp\else\^{}\fi}\catcode`\%=\active\def%{\%}$\mathdefault{10^{-11}}$}}%
\end{pgfscope}%
\begin{pgfscope}%
\pgfsetbuttcap%
\pgfsetroundjoin%
\definecolor{currentfill}{rgb}{0.000000,0.000000,0.000000}%
\pgfsetfillcolor{currentfill}%
\pgfsetlinewidth{0.501875pt}%
\definecolor{currentstroke}{rgb}{0.000000,0.000000,0.000000}%
\pgfsetstrokecolor{currentstroke}%
\pgfsetdash{}{0pt}%
\pgfsys@defobject{currentmarker}{\pgfqpoint{0.000000in}{0.000000in}}{\pgfqpoint{0.041667in}{0.000000in}}{%
\pgfpathmoveto{\pgfqpoint{0.000000in}{0.000000in}}%
\pgfpathlineto{\pgfqpoint{0.041667in}{0.000000in}}%
\pgfusepath{stroke,fill}%
}%
\begin{pgfscope}%
\pgfsys@transformshift{0.536040in}{1.383201in}%
\pgfsys@useobject{currentmarker}{}%
\end{pgfscope}%
\end{pgfscope}%
\begin{pgfscope}%
\definecolor{textcolor}{rgb}{0.000000,0.000000,0.000000}%
\pgfsetstrokecolor{textcolor}%
\pgfsetfillcolor{textcolor}%
\pgftext[x=0.248318in, y=1.340991in, left, base]{\color{textcolor}{\sffamily\fontsize{8.000000}{9.600000}\selectfont\catcode`\^=\active\def^{\ifmmode\sp\else\^{}\fi}\catcode`\%=\active\def%{\%}$\mathdefault{10^{-8}}$}}%
\end{pgfscope}%
\begin{pgfscope}%
\pgfsetbuttcap%
\pgfsetroundjoin%
\definecolor{currentfill}{rgb}{0.000000,0.000000,0.000000}%
\pgfsetfillcolor{currentfill}%
\pgfsetlinewidth{0.501875pt}%
\definecolor{currentstroke}{rgb}{0.000000,0.000000,0.000000}%
\pgfsetstrokecolor{currentstroke}%
\pgfsetdash{}{0pt}%
\pgfsys@defobject{currentmarker}{\pgfqpoint{0.000000in}{0.000000in}}{\pgfqpoint{0.041667in}{0.000000in}}{%
\pgfpathmoveto{\pgfqpoint{0.000000in}{0.000000in}}%
\pgfpathlineto{\pgfqpoint{0.041667in}{0.000000in}}%
\pgfusepath{stroke,fill}%
}%
\begin{pgfscope}%
\pgfsys@transformshift{0.536040in}{1.726711in}%
\pgfsys@useobject{currentmarker}{}%
\end{pgfscope}%
\end{pgfscope}%
\begin{pgfscope}%
\definecolor{textcolor}{rgb}{0.000000,0.000000,0.000000}%
\pgfsetstrokecolor{textcolor}%
\pgfsetfillcolor{textcolor}%
\pgftext[x=0.250484in, y=1.684502in, left, base]{\color{textcolor}{\sffamily\fontsize{8.000000}{9.600000}\selectfont\catcode`\^=\active\def^{\ifmmode\sp\else\^{}\fi}\catcode`\%=\active\def%{\%}$\mathdefault{10^{-5}}$}}%
\end{pgfscope}%
\begin{pgfscope}%
\pgfsetbuttcap%
\pgfsetroundjoin%
\definecolor{currentfill}{rgb}{0.000000,0.000000,0.000000}%
\pgfsetfillcolor{currentfill}%
\pgfsetlinewidth{0.501875pt}%
\definecolor{currentstroke}{rgb}{0.000000,0.000000,0.000000}%
\pgfsetstrokecolor{currentstroke}%
\pgfsetdash{}{0pt}%
\pgfsys@defobject{currentmarker}{\pgfqpoint{0.000000in}{0.000000in}}{\pgfqpoint{0.041667in}{0.000000in}}{%
\pgfpathmoveto{\pgfqpoint{0.000000in}{0.000000in}}%
\pgfpathlineto{\pgfqpoint{0.041667in}{0.000000in}}%
\pgfusepath{stroke,fill}%
}%
\begin{pgfscope}%
\pgfsys@transformshift{0.536040in}{2.070222in}%
\pgfsys@useobject{currentmarker}{}%
\end{pgfscope}%
\end{pgfscope}%
\begin{pgfscope}%
\definecolor{textcolor}{rgb}{0.000000,0.000000,0.000000}%
\pgfsetstrokecolor{textcolor}%
\pgfsetfillcolor{textcolor}%
\pgftext[x=0.248734in, y=2.028012in, left, base]{\color{textcolor}{\sffamily\fontsize{8.000000}{9.600000}\selectfont\catcode`\^=\active\def^{\ifmmode\sp\else\^{}\fi}\catcode`\%=\active\def%{\%}$\mathdefault{10^{-2}}$}}%
\end{pgfscope}%
\begin{pgfscope}%
\definecolor{textcolor}{rgb}{0.000000,0.000000,0.000000}%
\pgfsetstrokecolor{textcolor}%
\pgfsetfillcolor{textcolor}%
\pgftext[x=0.143096in,y=1.325949in,,bottom,rotate=90.000000]{\color{textcolor}{\sffamily\fontsize{9.000000}{10.800000}\selectfont\catcode`\^=\active\def^{\ifmmode\sp\else\^{}\fi}\catcode`\%=\active\def%{\%}Relative error in $\partial_x g_{tt}$}}%
\end{pgfscope}%
\begin{pgfscope}%
\pgfpathrectangle{\pgfqpoint{0.536040in}{0.352669in}}{\pgfqpoint{3.172850in}{1.946560in}}%
\pgfusepath{clip}%
\pgfsetbuttcap%
\pgfsetroundjoin%
\pgfsetlinewidth{1.204500pt}%
\definecolor{currentstroke}{rgb}{0.250980,0.400000,0.600000}%
\pgfsetstrokecolor{currentstroke}%
\pgfsetdash{{4.440000pt}{1.920000pt}}{0.000000pt}%
\pgfpathmoveto{\pgfqpoint{0.536040in}{0.541310in}}%
\pgfpathlineto{\pgfqpoint{3.708890in}{0.541310in}}%
\pgfusepath{stroke}%
\end{pgfscope}%
\begin{pgfscope}%
\pgfsetrectcap%
\pgfsetmiterjoin%
\pgfsetlinewidth{0.602250pt}%
\definecolor{currentstroke}{rgb}{0.000000,0.000000,0.000000}%
\pgfsetstrokecolor{currentstroke}%
\pgfsetdash{}{0pt}%
\pgfpathmoveto{\pgfqpoint{0.536040in}{0.352669in}}%
\pgfpathlineto{\pgfqpoint{0.536040in}{2.299229in}}%
\pgfusepath{stroke}%
\end{pgfscope}%
\begin{pgfscope}%
\pgfsetrectcap%
\pgfsetmiterjoin%
\pgfsetlinewidth{0.602250pt}%
\definecolor{currentstroke}{rgb}{0.000000,0.000000,0.000000}%
\pgfsetstrokecolor{currentstroke}%
\pgfsetdash{}{0pt}%
\pgfpathmoveto{\pgfqpoint{3.708890in}{0.352669in}}%
\pgfpathlineto{\pgfqpoint{3.708890in}{2.299229in}}%
\pgfusepath{stroke}%
\end{pgfscope}%
\begin{pgfscope}%
\pgfsetrectcap%
\pgfsetmiterjoin%
\pgfsetlinewidth{0.602250pt}%
\definecolor{currentstroke}{rgb}{0.000000,0.000000,0.000000}%
\pgfsetstrokecolor{currentstroke}%
\pgfsetdash{}{0pt}%
\pgfpathmoveto{\pgfqpoint{0.536040in}{0.352669in}}%
\pgfpathlineto{\pgfqpoint{3.708890in}{0.352669in}}%
\pgfusepath{stroke}%
\end{pgfscope}%
\begin{pgfscope}%
\pgfsetrectcap%
\pgfsetmiterjoin%
\pgfsetlinewidth{0.602250pt}%
\definecolor{currentstroke}{rgb}{0.000000,0.000000,0.000000}%
\pgfsetstrokecolor{currentstroke}%
\pgfsetdash{}{0pt}%
\pgfpathmoveto{\pgfqpoint{0.536040in}{2.299229in}}%
\pgfpathlineto{\pgfqpoint{3.708890in}{2.299229in}}%
\pgfusepath{stroke}%
\end{pgfscope}%
\begin{pgfscope}%
\pgfpathrectangle{\pgfqpoint{0.536040in}{0.352669in}}{\pgfqpoint{3.172850in}{1.946560in}}%
\pgfusepath{clip}%
\pgfsetrectcap%
\pgfsetroundjoin%
\pgfsetlinewidth{1.405250pt}%
\definecolor{currentstroke}{rgb}{0.768627,0.556863,0.211765}%
\pgfsetstrokecolor{currentstroke}%
\pgfsetdash{}{0pt}%
\pgfpathmoveto{\pgfqpoint{0.536040in}{2.252120in}}%
\pgfpathlineto{\pgfqpoint{0.542399in}{2.238517in}}%
\pgfpathlineto{\pgfqpoint{0.548757in}{2.221233in}}%
\pgfpathlineto{\pgfqpoint{0.555115in}{2.196830in}}%
\pgfpathlineto{\pgfqpoint{0.561474in}{2.152355in}}%
\pgfpathlineto{\pgfqpoint{0.567832in}{2.100022in}}%
\pgfpathlineto{\pgfqpoint{0.574191in}{2.175975in}}%
\pgfpathlineto{\pgfqpoint{0.580549in}{2.203247in}}%
\pgfpathlineto{\pgfqpoint{0.586907in}{2.219763in}}%
\pgfpathlineto{\pgfqpoint{0.593266in}{2.231401in}}%
\pgfpathlineto{\pgfqpoint{0.599624in}{2.240249in}}%
\pgfpathlineto{\pgfqpoint{0.605983in}{2.247289in}}%
\pgfpathlineto{\pgfqpoint{0.612341in}{2.181024in}}%
\pgfpathlineto{\pgfqpoint{0.618700in}{2.206033in}}%
\pgfpathlineto{\pgfqpoint{0.625058in}{2.221645in}}%
\pgfpathlineto{\pgfqpoint{0.631416in}{2.232798in}}%
\pgfpathlineto{\pgfqpoint{0.637775in}{2.241343in}}%
\pgfpathlineto{\pgfqpoint{0.644133in}{2.236556in}}%
\pgfpathlineto{\pgfqpoint{0.650492in}{2.218627in}}%
\pgfpathlineto{\pgfqpoint{0.656850in}{2.192801in}}%
\pgfpathlineto{\pgfqpoint{0.663208in}{2.142630in}}%
\pgfpathlineto{\pgfqpoint{0.669567in}{2.119341in}}%
\pgfpathlineto{\pgfqpoint{0.675925in}{2.180541in}}%
\pgfpathlineto{\pgfqpoint{0.682284in}{2.076887in}}%
\pgfpathlineto{\pgfqpoint{0.688642in}{2.172078in}}%
\pgfpathlineto{\pgfqpoint{0.695001in}{2.201184in}}%
\pgfpathlineto{\pgfqpoint{0.701359in}{2.204030in}}%
\pgfpathlineto{\pgfqpoint{0.707717in}{2.167482in}}%
\pgfpathlineto{\pgfqpoint{0.714076in}{1.950742in}}%
\pgfpathlineto{\pgfqpoint{0.720434in}{2.165341in}}%
\pgfpathlineto{\pgfqpoint{0.726793in}{2.142839in}}%
\pgfpathlineto{\pgfqpoint{0.733151in}{2.155253in}}%
\pgfpathlineto{\pgfqpoint{0.739509in}{2.091327in}}%
\pgfpathlineto{\pgfqpoint{0.745868in}{2.072448in}}%
\pgfpathlineto{\pgfqpoint{0.752226in}{2.073487in}}%
\pgfpathlineto{\pgfqpoint{0.758585in}{2.146378in}}%
\pgfpathlineto{\pgfqpoint{0.764943in}{2.130651in}}%
\pgfpathlineto{\pgfqpoint{0.771302in}{2.099913in}}%
\pgfpathlineto{\pgfqpoint{0.777660in}{2.150342in}}%
\pgfpathlineto{\pgfqpoint{0.784018in}{2.118726in}}%
\pgfpathlineto{\pgfqpoint{0.790377in}{2.141736in}}%
\pgfpathlineto{\pgfqpoint{0.796735in}{2.092652in}}%
\pgfpathlineto{\pgfqpoint{0.803094in}{2.088057in}}%
\pgfpathlineto{\pgfqpoint{0.809452in}{2.117537in}}%
\pgfpathlineto{\pgfqpoint{0.815810in}{2.068714in}}%
\pgfpathlineto{\pgfqpoint{0.822169in}{2.088087in}}%
\pgfpathlineto{\pgfqpoint{0.828527in}{2.111107in}}%
\pgfpathlineto{\pgfqpoint{0.834886in}{2.102592in}}%
\pgfpathlineto{\pgfqpoint{0.841244in}{2.113667in}}%
\pgfpathlineto{\pgfqpoint{0.847603in}{2.102583in}}%
\pgfpathlineto{\pgfqpoint{0.853961in}{1.999302in}}%
\pgfpathlineto{\pgfqpoint{0.866678in}{2.041688in}}%
\pgfpathlineto{\pgfqpoint{0.873036in}{2.066094in}}%
\pgfpathlineto{\pgfqpoint{0.879395in}{2.057855in}}%
\pgfpathlineto{\pgfqpoint{0.885753in}{2.098005in}}%
\pgfpathlineto{\pgfqpoint{0.892111in}{2.052354in}}%
\pgfpathlineto{\pgfqpoint{0.898470in}{2.050057in}}%
\pgfpathlineto{\pgfqpoint{0.904828in}{2.072491in}}%
\pgfpathlineto{\pgfqpoint{0.911187in}{2.037462in}}%
\pgfpathlineto{\pgfqpoint{0.917545in}{2.088005in}}%
\pgfpathlineto{\pgfqpoint{0.923904in}{2.004121in}}%
\pgfpathlineto{\pgfqpoint{0.930262in}{2.049763in}}%
\pgfpathlineto{\pgfqpoint{0.936620in}{1.970234in}}%
\pgfpathlineto{\pgfqpoint{0.942979in}{1.923670in}}%
\pgfpathlineto{\pgfqpoint{0.949337in}{2.011293in}}%
\pgfpathlineto{\pgfqpoint{0.955696in}{1.943141in}}%
\pgfpathlineto{\pgfqpoint{0.962054in}{1.957340in}}%
\pgfpathlineto{\pgfqpoint{0.968412in}{2.011152in}}%
\pgfpathlineto{\pgfqpoint{0.974771in}{1.898736in}}%
\pgfpathlineto{\pgfqpoint{0.981129in}{2.049055in}}%
\pgfpathlineto{\pgfqpoint{0.987488in}{2.050714in}}%
\pgfpathlineto{\pgfqpoint{0.993846in}{1.987224in}}%
\pgfpathlineto{\pgfqpoint{1.000205in}{1.997741in}}%
\pgfpathlineto{\pgfqpoint{1.006563in}{1.976320in}}%
\pgfpathlineto{\pgfqpoint{1.012921in}{2.020445in}}%
\pgfpathlineto{\pgfqpoint{1.019280in}{1.983975in}}%
\pgfpathlineto{\pgfqpoint{1.025638in}{1.990731in}}%
\pgfpathlineto{\pgfqpoint{1.031997in}{1.995726in}}%
\pgfpathlineto{\pgfqpoint{1.038355in}{1.968997in}}%
\pgfpathlineto{\pgfqpoint{1.044713in}{1.834945in}}%
\pgfpathlineto{\pgfqpoint{1.051072in}{2.013942in}}%
\pgfpathlineto{\pgfqpoint{1.057430in}{2.014744in}}%
\pgfpathlineto{\pgfqpoint{1.063789in}{1.961409in}}%
\pgfpathlineto{\pgfqpoint{1.070147in}{1.980318in}}%
\pgfpathlineto{\pgfqpoint{1.082864in}{1.947774in}}%
\pgfpathlineto{\pgfqpoint{1.089222in}{1.904078in}}%
\pgfpathlineto{\pgfqpoint{1.095581in}{1.899108in}}%
\pgfpathlineto{\pgfqpoint{1.101939in}{1.933174in}}%
\pgfpathlineto{\pgfqpoint{1.108298in}{1.985909in}}%
\pgfpathlineto{\pgfqpoint{1.114656in}{1.953903in}}%
\pgfpathlineto{\pgfqpoint{1.133731in}{1.973140in}}%
\pgfpathlineto{\pgfqpoint{1.140090in}{1.882378in}}%
\pgfpathlineto{\pgfqpoint{1.146448in}{1.957538in}}%
\pgfpathlineto{\pgfqpoint{1.152807in}{1.877829in}}%
\pgfpathlineto{\pgfqpoint{1.159165in}{1.923440in}}%
\pgfpathlineto{\pgfqpoint{1.165523in}{1.907824in}}%
\pgfpathlineto{\pgfqpoint{1.171882in}{1.923160in}}%
\pgfpathlineto{\pgfqpoint{1.178240in}{1.767488in}}%
\pgfpathlineto{\pgfqpoint{1.184599in}{1.848278in}}%
\pgfpathlineto{\pgfqpoint{1.190957in}{1.888518in}}%
\pgfpathlineto{\pgfqpoint{1.197315in}{1.871925in}}%
\pgfpathlineto{\pgfqpoint{1.203674in}{1.886495in}}%
\pgfpathlineto{\pgfqpoint{1.210032in}{1.880206in}}%
\pgfpathlineto{\pgfqpoint{1.216391in}{1.919824in}}%
\pgfpathlineto{\pgfqpoint{1.222749in}{1.817931in}}%
\pgfpathlineto{\pgfqpoint{1.229108in}{1.920793in}}%
\pgfpathlineto{\pgfqpoint{1.235466in}{1.876825in}}%
\pgfpathlineto{\pgfqpoint{1.241824in}{1.808433in}}%
\pgfpathlineto{\pgfqpoint{1.248183in}{1.891696in}}%
\pgfpathlineto{\pgfqpoint{1.254541in}{1.868210in}}%
\pgfpathlineto{\pgfqpoint{1.260900in}{1.768476in}}%
\pgfpathlineto{\pgfqpoint{1.267258in}{1.879796in}}%
\pgfpathlineto{\pgfqpoint{1.273616in}{1.861235in}}%
\pgfpathlineto{\pgfqpoint{1.279975in}{1.854831in}}%
\pgfpathlineto{\pgfqpoint{1.286333in}{1.887982in}}%
\pgfpathlineto{\pgfqpoint{1.292692in}{1.815176in}}%
\pgfpathlineto{\pgfqpoint{1.299050in}{1.839840in}}%
\pgfpathlineto{\pgfqpoint{1.305409in}{1.830201in}}%
\pgfpathlineto{\pgfqpoint{1.311767in}{1.829686in}}%
\pgfpathlineto{\pgfqpoint{1.318125in}{1.837849in}}%
\pgfpathlineto{\pgfqpoint{1.324484in}{1.779205in}}%
\pgfpathlineto{\pgfqpoint{1.330842in}{1.798907in}}%
\pgfpathlineto{\pgfqpoint{1.337201in}{1.864309in}}%
\pgfpathlineto{\pgfqpoint{1.343559in}{1.733274in}}%
\pgfpathlineto{\pgfqpoint{1.349917in}{1.858030in}}%
\pgfpathlineto{\pgfqpoint{1.356276in}{1.841713in}}%
\pgfpathlineto{\pgfqpoint{1.362634in}{1.781637in}}%
\pgfpathlineto{\pgfqpoint{1.368993in}{1.809746in}}%
\pgfpathlineto{\pgfqpoint{1.375351in}{1.735630in}}%
\pgfpathlineto{\pgfqpoint{1.381710in}{1.813915in}}%
\pgfpathlineto{\pgfqpoint{1.388068in}{1.752111in}}%
\pgfpathlineto{\pgfqpoint{1.394426in}{1.817436in}}%
\pgfpathlineto{\pgfqpoint{1.400785in}{1.824147in}}%
\pgfpathlineto{\pgfqpoint{1.407143in}{1.804636in}}%
\pgfpathlineto{\pgfqpoint{1.413502in}{1.799507in}}%
\pgfpathlineto{\pgfqpoint{1.419860in}{1.748566in}}%
\pgfpathlineto{\pgfqpoint{1.426218in}{1.812247in}}%
\pgfpathlineto{\pgfqpoint{1.432577in}{1.739991in}}%
\pgfpathlineto{\pgfqpoint{1.438935in}{1.777308in}}%
\pgfpathlineto{\pgfqpoint{1.445294in}{1.614586in}}%
\pgfpathlineto{\pgfqpoint{1.451652in}{1.776123in}}%
\pgfpathlineto{\pgfqpoint{1.458011in}{1.744839in}}%
\pgfpathlineto{\pgfqpoint{1.464369in}{1.726645in}}%
\pgfpathlineto{\pgfqpoint{1.470727in}{1.683303in}}%
\pgfpathlineto{\pgfqpoint{1.477086in}{1.697727in}}%
\pgfpathlineto{\pgfqpoint{1.483444in}{1.691647in}}%
\pgfpathlineto{\pgfqpoint{1.489803in}{1.736272in}}%
\pgfpathlineto{\pgfqpoint{1.496161in}{1.765661in}}%
\pgfpathlineto{\pgfqpoint{1.502519in}{1.692122in}}%
\pgfpathlineto{\pgfqpoint{1.508878in}{1.718971in}}%
\pgfpathlineto{\pgfqpoint{1.515236in}{1.763403in}}%
\pgfpathlineto{\pgfqpoint{1.521595in}{1.623796in}}%
\pgfpathlineto{\pgfqpoint{1.527953in}{1.724205in}}%
\pgfpathlineto{\pgfqpoint{1.534312in}{1.607750in}}%
\pgfpathlineto{\pgfqpoint{1.540670in}{1.703757in}}%
\pgfpathlineto{\pgfqpoint{1.547028in}{1.629296in}}%
\pgfpathlineto{\pgfqpoint{1.553387in}{1.732162in}}%
\pgfpathlineto{\pgfqpoint{1.559745in}{1.654607in}}%
\pgfpathlineto{\pgfqpoint{1.566104in}{1.429046in}}%
\pgfpathlineto{\pgfqpoint{1.572462in}{1.569109in}}%
\pgfpathlineto{\pgfqpoint{1.578820in}{1.657370in}}%
\pgfpathlineto{\pgfqpoint{1.585179in}{1.661796in}}%
\pgfpathlineto{\pgfqpoint{1.591537in}{1.696858in}}%
\pgfpathlineto{\pgfqpoint{1.597896in}{1.703420in}}%
\pgfpathlineto{\pgfqpoint{1.604254in}{1.690937in}}%
\pgfpathlineto{\pgfqpoint{1.610613in}{1.660974in}}%
\pgfpathlineto{\pgfqpoint{1.616971in}{1.669384in}}%
\pgfpathlineto{\pgfqpoint{1.623329in}{1.686913in}}%
\pgfpathlineto{\pgfqpoint{1.636046in}{1.553480in}}%
\pgfpathlineto{\pgfqpoint{1.642405in}{1.655875in}}%
\pgfpathlineto{\pgfqpoint{1.648763in}{1.657368in}}%
\pgfpathlineto{\pgfqpoint{1.655122in}{1.641121in}}%
\pgfpathlineto{\pgfqpoint{1.661480in}{1.682473in}}%
\pgfpathlineto{\pgfqpoint{1.667838in}{1.589618in}}%
\pgfpathlineto{\pgfqpoint{1.674197in}{1.473576in}}%
\pgfpathlineto{\pgfqpoint{1.680555in}{1.632547in}}%
\pgfpathlineto{\pgfqpoint{1.686914in}{1.662813in}}%
\pgfpathlineto{\pgfqpoint{1.693272in}{1.621993in}}%
\pgfpathlineto{\pgfqpoint{1.699630in}{1.526288in}}%
\pgfpathlineto{\pgfqpoint{1.705989in}{1.599662in}}%
\pgfpathlineto{\pgfqpoint{1.712347in}{1.611504in}}%
\pgfpathlineto{\pgfqpoint{1.718706in}{1.588327in}}%
\pgfpathlineto{\pgfqpoint{1.731423in}{1.650634in}}%
\pgfpathlineto{\pgfqpoint{1.737781in}{1.589519in}}%
\pgfpathlineto{\pgfqpoint{1.744139in}{1.602682in}}%
\pgfpathlineto{\pgfqpoint{1.750498in}{1.544293in}}%
\pgfpathlineto{\pgfqpoint{1.756856in}{1.560178in}}%
\pgfpathlineto{\pgfqpoint{1.763215in}{1.600056in}}%
\pgfpathlineto{\pgfqpoint{1.769573in}{1.609453in}}%
\pgfpathlineto{\pgfqpoint{1.775931in}{1.569249in}}%
\pgfpathlineto{\pgfqpoint{1.782290in}{1.547894in}}%
\pgfpathlineto{\pgfqpoint{1.788648in}{1.567637in}}%
\pgfpathlineto{\pgfqpoint{1.795007in}{1.591162in}}%
\pgfpathlineto{\pgfqpoint{1.801365in}{1.604674in}}%
\pgfpathlineto{\pgfqpoint{1.807724in}{1.538999in}}%
\pgfpathlineto{\pgfqpoint{1.814082in}{1.491346in}}%
\pgfpathlineto{\pgfqpoint{1.820440in}{1.555096in}}%
\pgfpathlineto{\pgfqpoint{1.826799in}{1.512354in}}%
\pgfpathlineto{\pgfqpoint{1.833157in}{1.563030in}}%
\pgfpathlineto{\pgfqpoint{1.839516in}{1.522551in}}%
\pgfpathlineto{\pgfqpoint{1.845874in}{1.519585in}}%
\pgfpathlineto{\pgfqpoint{1.852232in}{1.572972in}}%
\pgfpathlineto{\pgfqpoint{1.858591in}{1.584888in}}%
\pgfpathlineto{\pgfqpoint{1.864949in}{1.502750in}}%
\pgfpathlineto{\pgfqpoint{1.871308in}{1.432562in}}%
\pgfpathlineto{\pgfqpoint{1.877666in}{1.573810in}}%
\pgfpathlineto{\pgfqpoint{1.884025in}{1.531621in}}%
\pgfpathlineto{\pgfqpoint{1.890383in}{1.501566in}}%
\pgfpathlineto{\pgfqpoint{1.896741in}{1.501241in}}%
\pgfpathlineto{\pgfqpoint{1.903100in}{1.468185in}}%
\pgfpathlineto{\pgfqpoint{1.909458in}{1.537709in}}%
\pgfpathlineto{\pgfqpoint{1.915817in}{1.535375in}}%
\pgfpathlineto{\pgfqpoint{1.922175in}{1.554012in}}%
\pgfpathlineto{\pgfqpoint{1.928533in}{1.501022in}}%
\pgfpathlineto{\pgfqpoint{1.934892in}{1.542321in}}%
\pgfpathlineto{\pgfqpoint{1.941250in}{1.484825in}}%
\pgfpathlineto{\pgfqpoint{1.947609in}{1.524307in}}%
\pgfpathlineto{\pgfqpoint{1.953967in}{1.495705in}}%
\pgfpathlineto{\pgfqpoint{1.960326in}{1.462174in}}%
\pgfpathlineto{\pgfqpoint{1.973042in}{1.518411in}}%
\pgfpathlineto{\pgfqpoint{1.979401in}{1.494543in}}%
\pgfpathlineto{\pgfqpoint{1.985759in}{1.459854in}}%
\pgfpathlineto{\pgfqpoint{1.992118in}{1.444273in}}%
\pgfpathlineto{\pgfqpoint{1.998476in}{1.488640in}}%
\pgfpathlineto{\pgfqpoint{2.004834in}{1.507681in}}%
\pgfpathlineto{\pgfqpoint{2.011193in}{1.400797in}}%
\pgfpathlineto{\pgfqpoint{2.017551in}{1.435344in}}%
\pgfpathlineto{\pgfqpoint{2.023910in}{1.478555in}}%
\pgfpathlineto{\pgfqpoint{2.030268in}{1.387921in}}%
\pgfpathlineto{\pgfqpoint{2.036627in}{1.463151in}}%
\pgfpathlineto{\pgfqpoint{2.042985in}{1.259351in}}%
\pgfpathlineto{\pgfqpoint{2.049343in}{1.417193in}}%
\pgfpathlineto{\pgfqpoint{2.055702in}{1.406715in}}%
\pgfpathlineto{\pgfqpoint{2.062060in}{1.437367in}}%
\pgfpathlineto{\pgfqpoint{2.068419in}{1.412530in}}%
\pgfpathlineto{\pgfqpoint{2.074777in}{1.422325in}}%
\pgfpathlineto{\pgfqpoint{2.081135in}{1.454583in}}%
\pgfpathlineto{\pgfqpoint{2.087494in}{1.329250in}}%
\pgfpathlineto{\pgfqpoint{2.093852in}{1.402812in}}%
\pgfpathlineto{\pgfqpoint{2.100211in}{1.432207in}}%
\pgfpathlineto{\pgfqpoint{2.106569in}{1.427190in}}%
\pgfpathlineto{\pgfqpoint{2.112928in}{1.396116in}}%
\pgfpathlineto{\pgfqpoint{2.119286in}{1.319534in}}%
\pgfpathlineto{\pgfqpoint{2.125644in}{1.428290in}}%
\pgfpathlineto{\pgfqpoint{2.132003in}{1.357101in}}%
\pgfpathlineto{\pgfqpoint{2.138361in}{1.377157in}}%
\pgfpathlineto{\pgfqpoint{2.144720in}{1.303162in}}%
\pgfpathlineto{\pgfqpoint{2.151078in}{1.411560in}}%
\pgfpathlineto{\pgfqpoint{2.157436in}{1.346395in}}%
\pgfpathlineto{\pgfqpoint{2.163795in}{1.344551in}}%
\pgfpathlineto{\pgfqpoint{2.170153in}{1.240511in}}%
\pgfpathlineto{\pgfqpoint{2.176512in}{1.384900in}}%
\pgfpathlineto{\pgfqpoint{2.182870in}{1.393133in}}%
\pgfpathlineto{\pgfqpoint{2.189229in}{1.355426in}}%
\pgfpathlineto{\pgfqpoint{2.195587in}{1.303138in}}%
\pgfpathlineto{\pgfqpoint{2.201945in}{1.292367in}}%
\pgfpathlineto{\pgfqpoint{2.208304in}{1.278234in}}%
\pgfpathlineto{\pgfqpoint{2.214662in}{1.337684in}}%
\pgfpathlineto{\pgfqpoint{2.221021in}{1.382292in}}%
\pgfpathlineto{\pgfqpoint{2.227379in}{1.329133in}}%
\pgfpathlineto{\pgfqpoint{2.233737in}{1.264346in}}%
\pgfpathlineto{\pgfqpoint{2.240096in}{1.341702in}}%
\pgfpathlineto{\pgfqpoint{2.246454in}{1.298021in}}%
\pgfpathlineto{\pgfqpoint{2.252813in}{1.202547in}}%
\pgfpathlineto{\pgfqpoint{2.259171in}{1.193705in}}%
\pgfpathlineto{\pgfqpoint{2.265530in}{1.252495in}}%
\pgfpathlineto{\pgfqpoint{2.271888in}{1.224371in}}%
\pgfpathlineto{\pgfqpoint{2.278246in}{1.347127in}}%
\pgfpathlineto{\pgfqpoint{2.284605in}{1.268157in}}%
\pgfpathlineto{\pgfqpoint{2.290963in}{1.067987in}}%
\pgfpathlineto{\pgfqpoint{2.297322in}{1.169669in}}%
\pgfpathlineto{\pgfqpoint{2.303680in}{1.287922in}}%
\pgfpathlineto{\pgfqpoint{2.310038in}{1.265690in}}%
\pgfpathlineto{\pgfqpoint{2.316397in}{1.224614in}}%
\pgfpathlineto{\pgfqpoint{2.322755in}{1.282014in}}%
\pgfpathlineto{\pgfqpoint{2.329114in}{1.264105in}}%
\pgfpathlineto{\pgfqpoint{2.335472in}{1.200914in}}%
\pgfpathlineto{\pgfqpoint{2.341831in}{1.324569in}}%
\pgfpathlineto{\pgfqpoint{2.348189in}{1.230646in}}%
\pgfpathlineto{\pgfqpoint{2.354547in}{1.241815in}}%
\pgfpathlineto{\pgfqpoint{2.360906in}{1.259174in}}%
\pgfpathlineto{\pgfqpoint{2.367264in}{1.244246in}}%
\pgfpathlineto{\pgfqpoint{2.373623in}{1.246450in}}%
\pgfpathlineto{\pgfqpoint{2.386339in}{1.248565in}}%
\pgfpathlineto{\pgfqpoint{2.392698in}{1.195983in}}%
\pgfpathlineto{\pgfqpoint{2.399056in}{1.278713in}}%
\pgfpathlineto{\pgfqpoint{2.405415in}{1.245191in}}%
\pgfpathlineto{\pgfqpoint{2.411773in}{1.159138in}}%
\pgfpathlineto{\pgfqpoint{2.418132in}{1.204791in}}%
\pgfpathlineto{\pgfqpoint{2.424490in}{1.172438in}}%
\pgfpathlineto{\pgfqpoint{2.430848in}{1.159578in}}%
\pgfpathlineto{\pgfqpoint{2.437207in}{1.204886in}}%
\pgfpathlineto{\pgfqpoint{2.443565in}{1.175238in}}%
\pgfpathlineto{\pgfqpoint{2.449924in}{1.244539in}}%
\pgfpathlineto{\pgfqpoint{2.456282in}{1.149973in}}%
\pgfpathlineto{\pgfqpoint{2.462640in}{1.246849in}}%
\pgfpathlineto{\pgfqpoint{2.468999in}{1.137533in}}%
\pgfpathlineto{\pgfqpoint{2.475357in}{1.198771in}}%
\pgfpathlineto{\pgfqpoint{2.481716in}{1.229157in}}%
\pgfpathlineto{\pgfqpoint{2.488074in}{1.181860in}}%
\pgfpathlineto{\pgfqpoint{2.494433in}{1.155682in}}%
\pgfpathlineto{\pgfqpoint{2.500791in}{1.241794in}}%
\pgfpathlineto{\pgfqpoint{2.507149in}{1.119657in}}%
\pgfpathlineto{\pgfqpoint{2.513508in}{1.163913in}}%
\pgfpathlineto{\pgfqpoint{2.519866in}{1.122354in}}%
\pgfpathlineto{\pgfqpoint{2.526225in}{1.206521in}}%
\pgfpathlineto{\pgfqpoint{2.532583in}{1.098265in}}%
\pgfpathlineto{\pgfqpoint{2.538941in}{1.145399in}}%
\pgfpathlineto{\pgfqpoint{2.545300in}{1.084115in}}%
\pgfpathlineto{\pgfqpoint{2.551658in}{1.129023in}}%
\pgfpathlineto{\pgfqpoint{2.558017in}{1.179668in}}%
\pgfpathlineto{\pgfqpoint{2.564375in}{1.155295in}}%
\pgfpathlineto{\pgfqpoint{2.570734in}{1.143457in}}%
\pgfpathlineto{\pgfqpoint{2.577092in}{1.073499in}}%
\pgfpathlineto{\pgfqpoint{2.583450in}{1.025060in}}%
\pgfpathlineto{\pgfqpoint{2.589809in}{0.941458in}}%
\pgfpathlineto{\pgfqpoint{2.596167in}{1.154308in}}%
\pgfpathlineto{\pgfqpoint{2.602526in}{1.115561in}}%
\pgfpathlineto{\pgfqpoint{2.608884in}{1.134683in}}%
\pgfpathlineto{\pgfqpoint{2.615242in}{1.063734in}}%
\pgfpathlineto{\pgfqpoint{2.621601in}{1.094792in}}%
\pgfpathlineto{\pgfqpoint{2.627959in}{1.093506in}}%
\pgfpathlineto{\pgfqpoint{2.634318in}{1.146040in}}%
\pgfpathlineto{\pgfqpoint{2.640676in}{1.161181in}}%
\pgfpathlineto{\pgfqpoint{2.647035in}{1.102387in}}%
\pgfpathlineto{\pgfqpoint{2.653393in}{1.105905in}}%
\pgfpathlineto{\pgfqpoint{2.659751in}{1.038396in}}%
\pgfpathlineto{\pgfqpoint{2.666110in}{1.087409in}}%
\pgfpathlineto{\pgfqpoint{2.672468in}{0.988677in}}%
\pgfpathlineto{\pgfqpoint{2.678827in}{1.114461in}}%
\pgfpathlineto{\pgfqpoint{2.685185in}{1.073964in}}%
\pgfpathlineto{\pgfqpoint{2.691543in}{1.076081in}}%
\pgfpathlineto{\pgfqpoint{2.697902in}{1.024864in}}%
\pgfpathlineto{\pgfqpoint{2.704260in}{1.075703in}}%
\pgfpathlineto{\pgfqpoint{2.710619in}{1.051295in}}%
\pgfpathlineto{\pgfqpoint{2.716977in}{1.062329in}}%
\pgfpathlineto{\pgfqpoint{2.723336in}{1.003214in}}%
\pgfpathlineto{\pgfqpoint{2.729694in}{1.060194in}}%
\pgfpathlineto{\pgfqpoint{2.736052in}{1.083392in}}%
\pgfpathlineto{\pgfqpoint{2.742411in}{1.084517in}}%
\pgfpathlineto{\pgfqpoint{2.748769in}{1.127209in}}%
\pgfpathlineto{\pgfqpoint{2.755128in}{1.088627in}}%
\pgfpathlineto{\pgfqpoint{2.761486in}{1.132694in}}%
\pgfpathlineto{\pgfqpoint{2.767844in}{1.107054in}}%
\pgfpathlineto{\pgfqpoint{2.774203in}{1.117417in}}%
\pgfpathlineto{\pgfqpoint{2.780561in}{1.138811in}}%
\pgfpathlineto{\pgfqpoint{2.786920in}{1.149282in}}%
\pgfpathlineto{\pgfqpoint{2.793278in}{1.142051in}}%
\pgfpathlineto{\pgfqpoint{2.805995in}{1.165939in}}%
\pgfpathlineto{\pgfqpoint{2.812353in}{1.168952in}}%
\pgfpathlineto{\pgfqpoint{2.818712in}{1.179778in}}%
\pgfpathlineto{\pgfqpoint{2.825070in}{1.185000in}}%
\pgfpathlineto{\pgfqpoint{2.831429in}{1.187913in}}%
\pgfpathlineto{\pgfqpoint{2.844145in}{1.202814in}}%
\pgfpathlineto{\pgfqpoint{2.856862in}{1.215765in}}%
\pgfpathlineto{\pgfqpoint{2.882296in}{1.244099in}}%
\pgfpathlineto{\pgfqpoint{3.289235in}{1.684817in}}%
\pgfpathlineto{\pgfqpoint{3.677098in}{2.105745in}}%
\pgfpathlineto{\pgfqpoint{3.708890in}{2.141190in}}%
\pgfpathlineto{\pgfqpoint{3.708890in}{2.141190in}}%
\pgfusepath{stroke}%
\end{pgfscope}%
\begin{pgfscope}%
\definecolor{textcolor}{rgb}{0.431373,0.423529,0.462745}%
\pgfsetstrokecolor{textcolor}%
\pgfsetfillcolor{textcolor}%
\pgftext[x=3.220719in, y=1.762810in, left, base,rotate=38.000000]{\color{textcolor}{\sffamily\fontsize{7.000000}{8.400000}\selectfont\catcode`\^=\active\def^{\ifmmode\sp\else\^{}\fi}\catcode`\%=\active\def%{\%}$\propto h^2$}}%
\end{pgfscope}%
\begin{pgfscope}%
\definecolor{textcolor}{rgb}{0.431373,0.423529,0.462745}%
\pgfsetstrokecolor{textcolor}%
\pgfsetfillcolor{textcolor}%
\pgftext[x=1.875772in, y=1.833962in, left, base,rotate=340.000000]{\color{textcolor}{\sffamily\fontsize{7.000000}{8.400000}\selectfont\catcode`\^=\active\def^{\ifmmode\sp\else\^{}\fi}\catcode`\%=\active\def%{\%}$\propto h^{-1}$}}%
\end{pgfscope}%
\begin{pgfscope}%
\pgfsetroundcap%
\pgfsetroundjoin%
\pgfsetlinewidth{0.803000pt}%
\definecolor{currentstroke}{rgb}{0.768627,0.556863,0.211765}%
\pgfsetstrokecolor{currentstroke}%
\pgfsetdash{}{0pt}%
\pgfpathmoveto{\pgfqpoint{3.147444in}{1.171948in}}%
\pgfpathquadraticcurveto{\pgfqpoint{2.896572in}{1.140278in}}{\pgfqpoint{2.658024in}{1.110165in}}%
\pgfusepath{stroke}%
\end{pgfscope}%
\begin{pgfscope}%
\pgfsetroundcap%
\pgfsetroundjoin%
\pgfsetlinewidth{0.803000pt}%
\definecolor{currentstroke}{rgb}{0.768627,0.556863,0.211765}%
\pgfsetstrokecolor{currentstroke}%
\pgfsetdash{}{0pt}%
\pgfpathmoveto{\pgfqpoint{2.699042in}{1.095744in}}%
\pgfpathlineto{\pgfqpoint{2.658024in}{1.110165in}}%
\pgfpathlineto{\pgfqpoint{2.694172in}{1.134326in}}%
\pgfusepath{stroke}%
\end{pgfscope}%
\begin{pgfscope}%
\definecolor{textcolor}{rgb}{0.768627,0.556863,0.211765}%
\pgfsetstrokecolor{textcolor}%
\pgfsetfillcolor{textcolor}%
\pgftext[x=3.175242in,y=1.208826in,left,]{\color{textcolor}{\sffamily\fontsize{7.000000}{8.400000}\selectfont\catcode`\^=\active\def^{\ifmmode\sp\else\^{}\fi}\catcode`\%=\active\def%{\%}$\epsilon_{\mathrm{m}}^{2/3} \approx 10^{-11}$}}%
\end{pgfscope}%
\begin{pgfscope}%
\definecolor{textcolor}{rgb}{0.250980,0.400000,0.600000}%
\pgfsetstrokecolor{textcolor}%
\pgfsetfillcolor{textcolor}%
\pgftext[x=0.959087in,y=0.603608in,left,bottom]{\color{textcolor}{\sffamily\fontsize{7.000000}{8.400000}\selectfont\catcode`\^=\active\def^{\ifmmode\sp\else\^{}\fi}\catcode`\%=\active\def%{\%}$\sim\!\epsilon_{\mathrm{m}}$}}%
\end{pgfscope}%
\begin{pgfscope}%
\pgfsetroundcap%
\pgfsetroundjoin%
\pgfsetlinewidth{0.803000pt}%
\definecolor{currentstroke}{rgb}{0.431373,0.423529,0.462745}%
\pgfsetstrokecolor{currentstroke}%
\pgfsetdash{}{0pt}%
\pgfpathmoveto{\pgfqpoint{3.561042in}{1.063195in}}%
\pgfpathquadraticcurveto{\pgfqpoint{3.561042in}{0.822340in}}{\pgfqpoint{3.561042in}{0.581485in}}%
\pgfusepath{stroke}%
\end{pgfscope}%
\begin{pgfscope}%
\pgfsetroundcap%
\pgfsetroundjoin%
\pgfsetlinewidth{0.803000pt}%
\definecolor{currentstroke}{rgb}{0.431373,0.423529,0.462745}%
\pgfsetstrokecolor{currentstroke}%
\pgfsetdash{}{0pt}%
\pgfpathmoveto{\pgfqpoint{3.533264in}{1.007639in}}%
\pgfpathlineto{\pgfqpoint{3.561042in}{1.063195in}}%
\pgfpathlineto{\pgfqpoint{3.588819in}{1.007639in}}%
\pgfusepath{stroke}%
\end{pgfscope}%
\begin{pgfscope}%
\pgfsetroundcap%
\pgfsetroundjoin%
\pgfsetlinewidth{0.803000pt}%
\definecolor{currentstroke}{rgb}{0.431373,0.423529,0.462745}%
\pgfsetstrokecolor{currentstroke}%
\pgfsetdash{}{0pt}%
\pgfpathmoveto{\pgfqpoint{3.588819in}{0.637040in}}%
\pgfpathlineto{\pgfqpoint{3.561042in}{0.581485in}}%
\pgfpathlineto{\pgfqpoint{3.533264in}{0.637040in}}%
\pgfusepath{stroke}%
\end{pgfscope}%
\begin{pgfscope}%
\definecolor{textcolor}{rgb}{0.431373,0.423529,0.462745}%
\pgfsetstrokecolor{textcolor}%
\pgfsetfillcolor{textcolor}%
\pgftext[x=3.615038in,y=0.845152in,left,]{\color{textcolor}{\sffamily\fontsize{7.000000}{8.400000}\selectfont\catcode`\^=\active\def^{\ifmmode\sp\else\^{}\fi}\catcode`\%=\active\def%{\%}${\sim}\,10^5$}}%
\end{pgfscope}%
\begin{pgfscope}%
\pgfsetrectcap%
\pgfsetroundjoin%
\pgfsetlinewidth{1.405250pt}%
\definecolor{currentstroke}{rgb}{0.768627,0.556863,0.211765}%
\pgfsetstrokecolor{currentstroke}%
\pgfsetdash{}{0pt}%
\pgfpathmoveto{\pgfqpoint{0.613818in}{2.181613in}}%
\pgfpathlineto{\pgfqpoint{0.711040in}{2.181613in}}%
\pgfpathlineto{\pgfqpoint{0.808262in}{2.181613in}}%
\pgfusepath{stroke}%
\end{pgfscope}%
\begin{pgfscope}%
\definecolor{textcolor}{rgb}{0.000000,0.000000,0.000000}%
\pgfsetstrokecolor{textcolor}%
\pgfsetfillcolor{textcolor}%
\pgftext[x=0.886040in,y=2.147585in,left,base]{\color{textcolor}{\sffamily\fontsize{7.000000}{8.400000}\selectfont\catcode`\^=\active\def^{\ifmmode\sp\else\^{}\fi}\catcode`\%=\active\def%{\%}Central FD}}%
\end{pgfscope}%
\begin{pgfscope}%
\pgfsetbuttcap%
\pgfsetroundjoin%
\pgfsetlinewidth{1.204500pt}%
\definecolor{currentstroke}{rgb}{0.250980,0.400000,0.600000}%
\pgfsetstrokecolor{currentstroke}%
\pgfsetdash{{4.440000pt}{1.920000pt}}{0.000000pt}%
\pgfpathmoveto{\pgfqpoint{0.613818in}{2.038913in}}%
\pgfpathlineto{\pgfqpoint{0.711040in}{2.038913in}}%
\pgfpathlineto{\pgfqpoint{0.808262in}{2.038913in}}%
\pgfusepath{stroke}%
\end{pgfscope}%
\begin{pgfscope}%
\definecolor{textcolor}{rgb}{0.000000,0.000000,0.000000}%
\pgfsetstrokecolor{textcolor}%
\pgfsetfillcolor{textcolor}%
\pgftext[x=0.886040in,y=2.004885in,left,base]{\color{textcolor}{\sffamily\fontsize{7.000000}{8.400000}\selectfont\catcode`\^=\active\def^{\ifmmode\sp\else\^{}\fi}\catcode`\%=\active\def%{\%}AD (forward-mode)}}%
\end{pgfscope}%
\end{pgfpicture}%
\makeatother%
\endgroup%

\caption[Accuracy of finite differences vs.\ automatic differentiation]{%
  Relative error in $\pd{x} g_{tt}$ for the Schwarzschild
  Kerr--Schild metric at $r = 5M$, computed by central finite
  differences (gold) and forward-mode AD (blue dashed).  The
  finite-difference error follows the characteristic V-shape of
  \cref{eq:fd-total-error}: truncation error dominates for large
  step sizes ($E \propto h^2$, grey guide line), while round-off
  error dominates for small step sizes ($E \propto h^{-1}$).  The
  minimum near $h \approx 6 \times 10^{-6}$ reaches
  $\sim 10^{-11}$, consistent with the
  $\epsilon_{\text{m}}^{2/3}$ ceiling of \cref{eq:fd-optimal}.
  AD achieves $\sim 10^{-16}$ regardless of any step-size choice.}
\label{fig:ad-vs-fd-error}
\end{figure}

\paragraph{Computational cost.}

The accuracy advantage comes at comparable computational
cost.  An optimally implemented central-difference scheme
evaluates the metric at $x + h\,e_\mu$ and
$x - h\,e_\mu$ for each of the four coordinate directions,
giving $2 \times 4 = 8$ metric evaluations for all 40
derivatives.  AD requires four seeded evaluations, each
costing roughly twice a standard evaluation (every
floating-point operation is paired with a derivative
operation), giving $4 \times 2 = 8$ metric-equivalent
evaluations (\cref{sec:ad-forward}).  The two approaches
are roughly tied in cost, so accuracy is the
tiebreaker --- and AD wins by five orders of magnitude.
\implnote{The codebase's \texttt{christoffelsFiniteDiff}
recomputes the perturbed metric for each Christoffel
component rather than caching the eight perturbed
evaluations.  This is acceptable for a fallback path
that is rarely invoked.}

Analytic derivatives incur zero extra metric evaluations:
the derivative expressions are hardcoded functions of the
coordinates, evaluated directly.  For a production ray
tracer where every microsecond matters, this speed
advantage can be significant.  But the event-horizon
codebase is a reference implementation where correctness
takes priority over raw throughput, and the $8\times$
overhead of AD is a modest price for guaranteed accuracy.

\paragraph{Implementation effort.}

The three approaches reverse their ranking on the axis of
developer effort.  Analytic derivatives require
hand-deriving 40 expressions per spacetime --- products of
rational functions, chain-rule expansions through the
Kerr--Schild null vector, trigonometric terms for rotating
spacetimes.  A single sign error in one of the 40
components produces trajectories that silently violate the
Hamiltonian constraint.  The bug is difficult to isolate:
the incorrect derivative feeds into a contraction over
$\nu$ in the Christoffel formula, distributing the error
across multiple components of the equations of motion.

AD eliminates this class of errors.  The developer writes
the metric function once, polymorphically, and the library
produces all 40 derivatives mechanically.  The only
requirement is that the metric function uses the standard
\texttt{Floating} operations --- no hardcoded
\texttt{Double} constants, no calls to non-liftable
functions.  Every metric in the codebase satisfies this
constraint naturally.

Finite differences require no modification to the metric
function at all.  The fallback function
\texttt{christoffelsFiniteDiff} works with any
\texttt{Spacetime} instance by evaluating the existing
\texttt{metric4} method at perturbed points, making FD
the path of least resistance for spacetimes defined by
numerical data rather than closed-form expressions.

\paragraph{The practical verdict.}

AD occupies the sweet spot: it matches the accuracy of
correct analytic expressions at the cost of finite
differences, with implementation effort that scales with the
complexity of the metric itself rather than the complexity
of its derivatives.  \Cref{tab:ad-comparison} collects the
comparison.
%
\begin{table}[t]
\centering
\caption[Derivative computation strategies compared]{%
  Comparison of the three derivative computation strategies
  on accuracy, cost (metric evaluations for all 40
  derivatives), and implementation effort per spacetime.
  Both AD and analytic derivatives achieve accuracy limited
  only by floating-point round-off; the distinction is that
  analytic expressions are mathematically exact before
  evaluation, while AD accumulates round-off through the
  evaluation itself.}
\label{tab:ad-comparison}
\small
\begin{tabular}{@{}lccc@{}}
  \toprule
  & \textbf{AD} & \textbf{Finite diff.} & \textbf{Analytic} \\
  \midrule
  Accuracy
    & $\sim\!\epsilon_{\text{m}}$
    & $\sim\!\epsilon_{\text{m}}^{2/3}$
    & $\sim\!\epsilon_{\text{m}}$ \\[3pt]
  Metric evaluations
    & $\sim\!8$
    & $\sim\!8$
    & $0$ \\[3pt]
  Effort per spacetime
    & polymorphic $f$
    & none
    & 40 expressions \\
  \bottomrule
\end{tabular}
\end{table}
%
The codebase uses AD for every spacetime that has a
polymorphic metric function --- which is all of them ---
and reserves finite differences as a fallback for
hypothetical spacetimes defined by numerical data.
Analytic Christoffel symbols are not implemented for any
spacetime; the accuracy and maintainability gains from AD
make hand-coding unnecessary.
